\documentclass[11pt]{amsart}

\usepackage[T1]{fontenc}
\usepackage{lmodern}
\usepackage{microtype}
\usepackage{amsmath,amssymb,amsthm}
\usepackage[hidelinks]{hyperref}

\hypersetup{
  pdftitle={An Order-Three Counterexample to the Barat--Korondi--Varga Projection-Area Conjecture},
  pdfkeywords={cube dismantling, bootstrap percolation, orthogonal projection
area, independent set, grid graph, exhaustive enumeration}
}

\newtheorem{theorem}{Theorem}

\title[An order-three projection-area counterexample]
{An Order-Three Counterexample to the
Bar\'at--Korondi--Varga Projection-Area Conjecture}
\date{}

\subjclass[2020]{Primary 05C35; Secondary 05C69, 82B43}
\keywords{Cube dismantling, bootstrap percolation, orthogonal projection
area, independent set, grid graph, exhaustive enumeration}

\begin{document}

\begin{abstract}
Bar\'at, Korondi, and Varga conjectured that every solution to the
three-dimensional cube-dismantling problem has total orthogonal projection
area at least $n^2+6n-4$. We point out that the exhaustive order-three
catalogue of Bar\'at and Wanless already contains a counterexample. In
coordinates, one such solution has projection sizes $6,7,7$, and hence total
projection area $20<23$. We give a legal build-up certificate, and the same
catalogue yields the exact value $s_3(3)=20$. The configuration itself is not
new; the contribution is the resulting correction to the stated conjecture.
\end{abstract}

\maketitle

Write $[n]=\{1,\ldots,n\}$ and give $[n]^3$ its grid adjacency: two
vertices are adjacent if they differ by $1$ in exactly one coordinate. A
\emph{base position} is an independent set $B\subseteq[n]^3$ of size $n^2$.
It is a \emph{solution} if $[n]^3$ can be built from $B$ by adding the
vertices of $[n]^3\setminus B$ one at a time, each added vertex having
exactly three neighbours already present. This is the reverse of a legal
cube dismantling.

For $i\in\{1,2,3\}$, let $\pi_i$ delete the $i$th coordinate and set
\[
 A(B)=\sum_{i=1}^3 |\pi_i(B)|,
 \qquad
 s_3(n)=\min\{A(B):B\text{ is a solution in }[n]^3\}.
\]
In the notation above, a conjecture of Bar\'at, Korondi, and Varga
\cite{BKV}, stated as Conjecture~1 of Bar\'at and Wanless
\cite{BW},\footnote{The Bar\'at--Korondi--Varga formulation is quoted here at
second hand. Their extended abstract appeared in a symposium proceedings
volume that is not publicly available: no copy is indexed by arXiv, zbMATH,
Crossref, OpenAlex or Semantic Scholar, and the symposium's own abstract
repository is restricted to participants. We have therefore not seen their
statement, and in particular cannot confirm what number it carries there;
\cite{BW} attribute it to \cite{BKV} without one. The inequality displayed
here, including the quantifier range ``for all $n$'', is verbatim the
formulation printed as Conjecture~1 of \cite{BW}, and that formulation is
what the theorem below refutes.} asserts
\[
 s_3(n)\ge n^2+6n-4 \qquad\text{for all }n.
\]

\begin{theorem}
One has
\[
 s_3(3)=20<23=3^2+6\cdot3-4.
\]
Consequently, the Bar\'at--Korondi--Varga projection-area conjecture is
false as stated.
\end{theorem}

\begin{proof}
For compactness, write $ijk$ for $(i,j,k)$. Consider
\[
 B=\{111,113,122,131,223,311,322,331,333\}.
\]
No two vertices in $B$ are adjacent, so $B$ is a base position. Its
projections are
\begin{align*}
 \pi_1(B)&=\{(1,1),(1,3),(2,2),(2,3),(3,1),(3,3)\},\\
 \pi_2(B)&=\{(1,1),(1,2),(1,3),(2,3),(3,1),(3,2),(3,3)\},\\
 \pi_3(B)&=\{(1,1),(1,2),(1,3),(2,2),(3,1),(3,2),(3,3)\}.
\end{align*}
Thus $A(B)=6+7+7=20$.

It remains to certify that $B$ is a solution. Put $S_0=B$, and for a set
$S$ let $N_S(v)$ be the set of neighbours of $v$ in $S$. In the following
table, set $S_k=S_{k-1}\cup\{v_k\}$. Each displayed set is the full current
neighbourhood of the added vertex.

\begin{center}
\small
\renewcommand{\arraystretch}{1.06}
\begin{tabular}{c@{\;}c@{\;}l@{\qquad}c@{\;}c@{\;}l}
$k$ & $v_k$ & $N_{S_{k-1}}(v_k)$
& $k$ & $v_k$ & $N_{S_{k-1}}(v_k)$\\
\hline
 1 & 112 & $\{111,113,122\}$
& 10 & 232 & $\{222,231,332\}$\\
 2 & 121 & $\{111,122,131\}$
& 11 & 233 & $\{223,232,333\}$\\
 3 & 123 & $\{113,122,223\}$
& 12 & 132 & $\{122,131,232\}$\\
 4 & 222 & $\{122,223,322\}$
& 13 & 133 & $\{123,132,233\}$\\
 5 & 321 & $\{311,322,331\}$
& 14 & 211 & $\{111,221,311\}$\\
 6 & 323 & $\{223,322,333\}$
& 15 & 212 & $\{112,211,222\}$\\
 7 & 332 & $\{322,331,333\}$
& 16 & 213 & $\{113,212,223\}$\\
 8 & 221 & $\{121,222,321\}$
& 17 & 312 & $\{212,311,322\}$\\
 9 & 231 & $\{131,221,331\}$
& 18 & 313 & $\{213,312,323\}$
\end{tabular}
\end{center}

Every current neighbourhood has cardinality three, and
$\{v_1,\ldots,v_{18}\}=[3]^3\setminus B$. Hence the table is a legal
build-up from $B$ to $[3]^3$, proving $s_3(3)\le20$.

For the reverse inequality, Bar\'at and Wanless exhaustively enumerated all
order-three solutions and found seven isometry classes, one representative of
each being drawn in their Figure~6 \cite[Section~4, Table~1 and
Figure~6]{BW}. Computing the projection totals of those seven representatives
gives
\[
 27,\ 27,\ 25,\ 23,\ 21,\ 20,\ 21,
\]
of which only the multiset is used here. A cube isometry preserves total
projection area, so these are the totals of the seven classes. Their
classification therefore implies $s_3(3)\ge20$, and hence $s_3(3)=20$. The
same conclusion is reached without the figure, by the exhaustive enumeration
recorded at the end of this note: of the $4224$ independent $9$-subsets of
$[3]^3$, exactly $116$ are solutions, they fall into seven isometry classes
carrying the same seven totals, and the least of those totals is $20$.
\end{proof}

\paragraph{Prior work.}
The set $B$ is isometric to one of the seven representatives in Figure~6 of
\cite{BW}, necessarily the one of total projection area $20$, since exactly
one of the seven classes attains that value. Thus the configuration, its
validity as a solution, and the classification that pins $s_3(3)$ were all
already contained in the published order-three catalogue; the observation
recorded here is that the resulting value contradicts the conjectured bound.

\paragraph{Status and scope.}
What is refuted is the formulation displayed above, namely the one printed as
Conjecture~1 of \cite{BW}; on the sourcing of that wording, and on what can
and cannot be checked about the numbering it carries in \cite{BKV}, see the
footnote there. The failure is not confined to the minimiser: three of the
seven order-three isometry classes have total projection area below $23$,
namely $20$, $21$ and $21$. Nor is the bound false at every order. It holds
with equality at the two smaller ones, $s_3(1)=3$ and $s_3(2)=12$, so $n=3$
is the first order at which it fails. No claim is made about $s_3(n)$ for
$n\ge4$.

\paragraph{Exact verification.}
The computation reported above was re-run independently, on a separate
machine, by a standalone program that takes the data printed here as input
and recomputes everything else from the definitions in exact integer
arithmetic. All $32$ of its checks agree with what is stated here. It
rebuilds the grid graph on $[3]^3$ and the addition rule, and separates the
rule ``exactly three neighbours already present'' from ``at least three'' by
exhibiting configurations on which the two readings differ; it re-derives
the independence of $B$, the three projection sets elementwise, and
$A(B)=6+7+7=20$; it replays all eighteen rows of the table, requiring each
recomputed neighbourhood to equal the printed one and to have exactly three
elements, and it finds a legal addition order of its own without consulting
the table; and it carries out the order-three census afresh, enumerating all
$4224$ independent $9$-subsets of $[3]^3$, testing each for solution-hood,
building and verifying the $48$-element cube isometry group, checking that
$A$ is constant on isometry classes, and obtaining $116$ solutions in seven
classes with least total $20$. The count $4224$ is confirmed by a second,
unrelated algorithm, a layer transfer-matrix dynamic program; the same
census at $n=1$ and $n=2$ returns equality in the conjectured bound, a
control against systematic undershoot; and a mutation test reports failure
under each of a suite of deliberate corruptions of the printed data and of
the program's own geometry, with one exception that is provably immaterial
on this input, namely interchanging the first two rows of the table, whose
printed neighbourhoods both lie inside $B$. No file accompanies this note:
every quantity above is recomputable from the data printed here together
with the definitions of the first two paragraphs, and the program used and
its transcript are retained in an auxiliary archive, available on request.

Two things the re-run did not reproduce. First, nothing beyond $n=3$: a
census at $n=4$ would have to test $15{,}183{,}071{,}352$ independent
$16$-subsets of $[4]^3$, about $3.6$ million times the order-three workload,
and it was not attempted---nor is anything beyond $n=3$ claimed here.
Second, the attributions, which are assertions about the literature rather
than about mathematics: that the inequality transcribed above is the one
posed in \cite{BKV}, and that the seven representatives drawn in Figure~6 of
\cite{BW} are those of the seven isometry classes whose totals are computed
above---only the multiset of those totals is used, which is all the theorem
needs. Those are readings of the cited papers, and in the case of \cite{BKV}
of a source we were unable to consult at all. The re-run agrees with every
number printed here; agreement is evidence that the computation is right,
not a proof of it.

\begin{thebibliography}{9}

\bibitem{BKV}
J.~Bar\'at, M.~Korondi, and V.~Varga,
\emph{How to dismantle an atomic cube in zero gravity},
in \emph{Proceedings of the 7th Hungarian--Japanese Symposium on Discrete
Mathematics and Its Applications}, Kyoto, Japan, 2011; extended abstract, not
publicly available. The details of this entry are those printed in the
reference list of \cite{BW}; we have not seen the volume.

\bibitem{BW}
J.~Bar\'at and I.~M.~Wanless,
\emph{A cube dismantling problem related to bootstrap percolation},
J. Stat. Phys. \textbf{149} (2012), 754--770,
\href{https://doi.org/10.1007/s10955-012-0622-7}
{doi:10.1007/s10955-012-0622-7};
\href{https://arxiv.org/abs/2603.23913}{arXiv:2603.23913} [math.CO].

\end{thebibliography}

\end{document}